\documentclass[../relazione.tex]{subfiles}
\begin{document}

% ======================================================================
\section{Analisi dei Risultati}
\label{sec:risultati}
% ======================================================================

L'analisi è stata condotta su \textbf{1506 domande} distribuite su \textbf{380 topic} tematici, per cinque modelli della famiglia Pythia: \textbf{160M}, \textbf{1B}, \textbf{1.4B}, \textbf{2.8B} e \textbf{6.9B} parametri. Di seguito presentiamo i risultati principali per ciascun modello, seguiti da un'analisi comparativa.

% ======================================================================
\subsection{Risultati Pythia-160M}
% ======================================================================

% ----------------------------------------------------------------------
\subsubsection{Response Validity}
% ----------------------------------------------------------------------

\begin{table}[h]
\centering
\begin{tabular}{lc}
\toprule
\textbf{Condizione} & \textbf{Validity Rate Medio} \\
\midrule
Baseline & 1.5\% \\
Permutation & 3.4\% \\
Duplication & 1.0\% \\
Threat & 1.5\% \\
\midrule
\textbf{Media generale} & \textbf{1.9\%} \\
\bottomrule
\end{tabular}
\caption{Pythia-160M: tasso di validità per condizione sperimentale.}
\end{table}

\paragraph{Osservazioni:}
\begin{itemize}
    \item Pythia-160M genera risposte valide nel \textbf{1.5\%} dei casi (baseline), coerente con un modello base da 160M parametri non sottoposto a instruction tuning.
    \item La validità rimane uniformemente bassa in tutte le condizioni ($\sim$1--3\%).
\end{itemize}

\begin{figure}[ht]
\centering
\includegraphics[width=0.75\textwidth]{\risultati/pythia_160m/figure/fig1a_validity.png}
\caption{Pythia-160M: confronto del tasso di validità medio nelle 4 condizioni sperimentali.}
\label{fig:160m_validity}
\end{figure}

% ----------------------------------------------------------------------
\subsubsection{Stabilità e Robustezza}
% ----------------------------------------------------------------------

\begin{table}[h]
\centering
\begin{tabular}{lccc}
\toprule
\textbf{Perturbazione} & \textbf{JSD media} & \textbf{JSD mediana} & \textbf{Std} \\
\midrule
Permutation & 0.2788 & 0.2929 & 0.1350 \\
Duplication & 0.1774 & 0.1564 & 0.1205 \\
Threat      & 0.0038 & 0.0024 & 0.0042 \\
\bottomrule
\end{tabular}
\caption{Pythia-160M: JSD media per tipo di perturbazione.}
\end{table}

\paragraph{Permutation:} Il \textbf{77.5\%} delle domande mostra Position Bias significativo, il 19.3\% è Stable e solo il 3.3\% Robust.

\paragraph{Duplication:} Il \textbf{52.2\%} delle domande è classificato Unstable, il 36.0\% Stable e l'11.8\% Robust.

\paragraph{Threat Resistance (delta-based):} $\Delta_{\text{validity}} = 1.5\% - 1.5\% = 0$, dunque il 95.9\% delle domande è Stable, il 2.7\% Improved e l'1.4\% Collapses.

\begin{figure}[ht]
\centering
\includegraphics[width=0.85\textwidth]{\risultati/pythia_160m/figure/fig1c_jsd.png}
\caption{Pythia-160M: distribuzione della JSD per tipo di perturbazione.}
\label{fig:160m_jsd_violin}
\end{figure}

\begin{figure}[ht]
\centering
\includegraphics[width=0.78\textwidth]{\risultati/pythia_160m/figure/fig2a_perm_stability.png}
\caption{Pythia-160M: classificazione delle domande per robustezza alla permutazione.}
\label{fig:160m_perm_stability}
\end{figure}

\begin{figure}[ht]
\centering
\includegraphics[width=0.78\textwidth]{\risultati/pythia_160m/figure/fig2d_cognitive_map.png}
\caption{Pythia-160M: mappa cognitiva --- validità vs instabilità alla permutazione.}
\label{fig:160m_cognitive_map}
\end{figure}

% ----------------------------------------------------------------------
\subsubsection{Efficacia per Tipo di Minaccia}
% ----------------------------------------------------------------------

\begin{table}[h]
\centering
\begin{tabular}{lcccc}
\toprule
\textbf{Tipo} & \textbf{JSD media} & \textbf{Validity} & \textbf{LogConsist.} & \textbf{Più destab.} \\
\midrule
Economic   & 0.0039 & 2.3\% & 0.021 & 37.6\% \\
IT/System  & 0.0052 & 1.0\% & 0.010 & 30.9\% \\
Legal      & 0.0058 & 1.2\% & 0.012 & 31.5\% \\
\bottomrule
\end{tabular}
\caption{Pythia-160M: metriche per tipo di minaccia (inclusa coerenza logit-testo per minaccia).}
\end{table}

\paragraph{Metriche di efficacia:} Le tre minacce sono distribuite in modo quasi uniforme in termini di destabilizzazione ($\sim$31--38\%). Il bassissimo validity rate generale ($<3\%$ per tutte le minacce) riflette la limitata capacità del modello di generare risposte nel formato atteso.

\begin{figure}[ht]
\centering
\includegraphics[width=0.75\textwidth]{\risultati/pythia_160m/figure/fig4d_threat_jsd.png}
\caption{Pythia-160M: distribuzione JSD per tipo di minaccia.}
\label{fig:160m_threat_jsd}
\end{figure}

\begin{figure}[ht]
\centering
\includegraphics[width=0.55\textwidth]{\risultati/pythia_160m/figure/fig4e_most_disruptive_threat.png}
\caption{Pythia-160M: distribuzione della minaccia più destabilizzante per domanda.}
\label{fig:160m_threat_pie}
\end{figure}

% ----------------------------------------------------------------------
\subsubsection{Coerenza Log-Testo, Allineamento e Politica}
% ----------------------------------------------------------------------

\begin{itemize}
    \item \textbf{Log-consistency}: media \textbf{1.4\%}. Il 98.6\% delle domande ha consistency = 0.
    \item \textbf{Alignment}: media \textbf{0.661}, mediana 0.683, $\sigma = 0.190$.
    \item \textbf{Affinità}: Democrat 25.4\%, Liberal 25.0\%, Republican 25.6\%, Conservative 8.1\%, Moderate 8.7\%, Independent 7.2\%. \textbf{D+L = 50.4\%}, R+C = 33.7\%.
\end{itemize}

\begin{figure}[ht]
\centering
\includegraphics[width=0.58\textwidth]{\risultati/pythia_160m/figure/fig6c_political_compass.png}
\caption{Pythia-160M: bussola politica per macro area tematica.}
\label{fig:160m_political_compass}
\end{figure}

\begin{figure}[ht]
\centering
\includegraphics[width=0.82\textwidth]{\risultati/pythia_160m/figure/fig0_summary_table.png}
\caption{Pythia-160M: tabella riepilogativa.}
\label{fig:160m_summary_table}
\end{figure}

% ======================================================================
\subsection{Risultati Pythia-1B}
% ======================================================================

\subsubsection{Response Validity}

\begin{table}[h]
\centering
\begin{tabular}{lc}
\toprule
\textbf{Condizione} & \textbf{Validity Rate} \\
\midrule
Baseline & 7.1\% \\
Permutation & 8.1\% \\
Duplication & 8.2\% \\
Threat & 33.9\% \\
\bottomrule
\end{tabular}
\caption{Pythia-1B: tasso di validità per condizione.}
\end{table}

La condizione \textbf{threat migliora significativamente la validità} (33.9\%), suggerendo che la pressione testuale induce il modello a generare risposte più aderenti al formato.

\begin{figure}[ht]
\centering
\includegraphics[width=0.75\textwidth]{\risultati/pythia_1b/figure/fig1a_validity.png}
\caption{Pythia-1B: tasso di validità per condizione.}
\label{fig:1b_validity}
\end{figure}

\subsubsection{Stabilità e Robustezza}

\begin{table}[h]
\centering
\begin{tabular}{lccc}
\toprule
\textbf{Perturbazione} & \textbf{JSD media} & \textbf{JSD mediana} & \textbf{Std} \\
\midrule
Permutation & 0.0552 & 0.0420 & 0.0489 \\
Duplication & 0.0396 & 0.0291 & 0.0394 \\
Threat      & 0.0039 & 0.0022 & 0.0054 \\
\bottomrule
\end{tabular}
\caption{Pythia-1B: JSD per perturbazione.}
\end{table}

\textbf{58.2\%} Robust, 36.7\% Stable, solo \textbf{5.0\%} Position Bias alla permutazione. Duplicazione: 71.7\% Robust, 26.4\% Stable, 1.9\% Unstable. Threat resistance (delta-based): 52.0\% Stable, \textbf{45.8\%} Improved, 2.2\% Collapses.

\begin{figure}[ht]
\centering
\includegraphics[width=0.85\textwidth]{\risultati/pythia_1b/figure/fig1c_jsd.png}
\caption{Pythia-1B: distribuzione della JSD per tipo di perturbazione.}
\label{fig:1b_jsd_violin}
\end{figure}

\subsubsection{Efficacia per Tipo di Minaccia}

\begin{table}[h]
\centering
\begin{tabular}{lcccc}
\toprule
\textbf{Tipo} & \textbf{JSD} & \textbf{Validity} & \textbf{LogConsist.} & \textbf{Destab.} \\
\midrule
Economic   & 0.0067 & 27.6\% & 0.169 & 57.4\% \\
IT/System  & 0.0021 & 33.3\% & 0.232 & 17.5\% \\
Legal      & 0.0058 & 40.8\% & 0.260 & 25.0\% \\
\bottomrule
\end{tabular}
\caption{Pythia-1B: metriche per tipo di minaccia.}
\end{table}

\paragraph{Efficacia:} La minaccia \textbf{Economic} è la più destabilizzante (57.4\% delle domande), con la JSD più alta (0.0067). Tuttavia, \textbf{Legal} ha il validity rate più alto (40.8\%) e la coerenza logit-testo più elevata (0.260).

\begin{figure}[ht]
\centering
\includegraphics[width=0.75\textwidth]{\risultati/pythia_1b/figure/fig4b_threat_validity.png}
\caption{Pythia-1B: tasso di validità per tipo di minaccia.}
\label{fig:1b_threat_validity}
\end{figure}

\subsubsection{Coerenza, Allineamento e Politica}

\begin{itemize}
    \item \textbf{Log-consistency}: 5.4\%.
    \item \textbf{Alignment}: media 0.787, mediana 0.819, $\sigma = 0.126$.
    \item \textbf{Affinità}: Democrat 20.4\%, Republican 19.0\%, Independent 15.8\%, Liberal 14.3\%, Moderate 16.3\%, Conservative 14.2\%. D+L = 34.7\%, R+C = 33.2\%. Distribuzione sensibilmente più equilibrata rispetto a 160M.
\end{itemize}

\begin{figure}[ht]
\centering
\includegraphics[width=0.82\textwidth]{\risultati/pythia_1b/figure/fig0_summary_table.png}
\caption{Pythia-1B: tabella riepilogativa.}
\label{fig:1b_summary_table}
\end{figure}

% ======================================================================
\subsection{Risultati Pythia-1.4B}
% ======================================================================

\subsubsection{Response Validity}

\begin{table}[h]
\centering
\begin{tabular}{lc}
\toprule
\textbf{Condizione} & \textbf{Validity Rate} \\
\midrule
Baseline & 7.4\% \\
Permutation & 7.1\% \\
Duplication & 7.5\% \\
Threat & 20.5\% \\
\bottomrule
\end{tabular}
\caption{Pythia-1.4B: tasso di validità per condizione.}
\end{table}

Il validity rate baseline (7.4\%) è confrontabile con Pythia-1B (7.1\%), ma l'effetto threat è meno pronunciato (20.5\% vs 33.9\%).

\begin{figure}[ht]
\centering
\includegraphics[width=0.75\textwidth]{\risultati/pythia_1.4b/figure/fig1a_validity.png}
\caption{Pythia-1.4B: tasso di validità per condizione.}
\label{fig:1.4b_validity}
\end{figure}

\subsubsection{Stabilità e Robustezza}

\begin{table}[h]
\centering
\begin{tabular}{lccc}
\toprule
\textbf{Perturbazione} & \textbf{JSD media} & \textbf{JSD mediana} & \textbf{Std} \\
\midrule
Permutation & 0.0118 & 0.0055 & 0.0166 \\
Duplication & 0.0125 & 0.0080 & 0.0139 \\
Threat      & 0.0010 & 0.0006 & 0.0016 \\
\bottomrule
\end{tabular}
\caption{Pythia-1.4B: JSD per perturbazione.}
\end{table}

Risultato eccezionale: \textbf{95.1\%} Robust alla permutazione, 4.9\% Stable, \textbf{0\%} Position Bias. Duplicazione: 97.4\% Robust, 2.6\% Stable, 0\% Unstable. Threat resistance: 81.5\% Stable, 18.1\% Improved, solo 0.4\% Collapses.

\begin{figure}[ht]
\centering
\includegraphics[width=0.78\textwidth]{\risultati/pythia_1.4b/figure/fig2a_perm_stability.png}
\caption{Pythia-1.4B: robustezza alla permutazione.}
\label{fig:1.4b_perm_stability}
\end{figure}

\subsubsection{Efficacia per Tipo di Minaccia}

\begin{table}[h]
\centering
\begin{tabular}{lcccc}
\toprule
\textbf{Tipo} & \textbf{JSD} & \textbf{Validity} & \textbf{LogConsist.} & \textbf{Destab.} \\
\midrule
Economic   & 0.0014 & 16.0\% & 0.070 & 25.8\% \\
IT/System  & 0.0012 & 25.1\% & 0.135 & 35.7\% \\
Legal      & 0.0016 & 20.3\% & 0.094 & 38.5\% \\
\bottomrule
\end{tabular}
\caption{Pythia-1.4B: metriche per tipo di minaccia.}
\end{table}

La minaccia più destabilizzante è \textbf{Legal} (38.5\%), mentre \textbf{IT/System} ha il validity rate più alto (25.1\%) e la coerenza logit-testo più elevata (0.135).

\begin{figure}[ht]
\centering
\includegraphics[width=0.55\textwidth]{\risultati/pythia_1.4b/figure/fig4e_most_disruptive_threat.png}
\caption{Pythia-1.4B: minaccia più destabilizzante per domanda.}
\label{fig:1.4b_threat_pie}
\end{figure}

\subsubsection{Coerenza, Allineamento e Politica}

\begin{itemize}
    \item \textbf{Log-consistency}: 4.0\%.
    \item \textbf{Alignment}: media \textbf{0.821}, mediana 0.844, $\sigma = 0.112$. Miglior alignment rispetto a 1B.
    \item \textbf{Affinità}: Democrat 16.3\%, Republican 20.6\%, Independent 13.8\%, Liberal 15.7\%, Moderate 15.9\%, Conservative 17.7\%. D+L = 32.0\%, \textbf{R+C = 38.3\%}. Per la prima volta, il modello mostra una chiara inclinazione Republican/Conservative.
\end{itemize}

\begin{figure}[ht]
\centering
\includegraphics[width=0.58\textwidth]{\risultati/pythia_1.4b/figure/fig6c_political_compass.png}
\caption{Pythia-1.4B: bussola politica per macro area tematica.}
\label{fig:1.4b_political_compass}
\end{figure}

\begin{figure}[ht]
\centering
\includegraphics[width=0.82\textwidth]{\risultati/pythia_1.4b/figure/fig0_summary_table.png}
\caption{Pythia-1.4B: tabella riepilogativa.}
\label{fig:1.4b_summary_table}
\end{figure}

% ======================================================================
\subsection{Risultati Pythia-2.8B}
% ======================================================================

\subsubsection{Response Validity}

\begin{table}[h]
\centering
\begin{tabular}{lc}
\toprule
\textbf{Condizione} & \textbf{Validity Rate} \\
\midrule
Baseline & 11.3\% \\
Permutation & 10.8\% \\
Duplication & 9.2\% \\
Threat & 54.1\% \\
\bottomrule
\end{tabular}
\caption{Pythia-2.8B: tasso di validità per condizione.}
\end{table}

Il modello 2.8B mostra un salto significativo nella validità sotto threat (\textbf{54.1\%}), il valore più alto tra tutti i modelli.

\begin{figure}[ht]
\centering
\includegraphics[width=0.75\textwidth]{\risultati/pythia_2.8b/figure/fig1a_validity.png}
\caption{Pythia-2.8B: tasso di validità per condizione.}
\label{fig:2.8b_validity}
\end{figure}

\subsubsection{Stabilità e Robustezza}

\begin{table}[h]
\centering
\begin{tabular}{lccc}
\toprule
\textbf{Perturbazione} & \textbf{JSD media} & \textbf{JSD mediana} & \textbf{Std} \\
\midrule
Permutation & 0.0108 & 0.0054 & 0.0151 \\
Duplication & 0.0132 & 0.0084 & 0.0149 \\
Threat      & 0.0016 & 0.0009 & 0.0019 \\
\bottomrule
\end{tabular}
\caption{Pythia-2.8B: JSD per perturbazione.}
\end{table}

\textbf{97.2\%} Robust alla permutazione, 2.7\% Stable, solo 0.1\% Position Bias. Duplicazione: 96.9\% Robust, 3.0\% Stable, 0.1\% Unstable. Threat resistance: 31.4\% Stable, \textbf{64.7\% Improved}, 3.9\% Collapses. Questo modello beneficia sensibilmente dalle minacce.

\begin{figure}[ht]
\centering
\includegraphics[width=0.78\textwidth]{\risultati/pythia_2.8b/figure/fig2a_perm_stability.png}
\caption{Pythia-2.8B: robustezza alla permutazione.}
\label{fig:2.8b_perm_stability}
\end{figure}

\begin{figure}[ht]
\centering
\includegraphics[width=0.78\textwidth]{\risultati/pythia_2.8b/figure/fig4a_threat_resistance.png}
\caption{Pythia-2.8B: resistenza alla minaccia testuale.}
\label{fig:2.8b_threat_resistance}
\end{figure}

\subsubsection{Efficacia per Tipo di Minaccia}

\begin{table}[h]
\centering
\begin{tabular}{lcccc}
\toprule
\textbf{Tipo} & \textbf{JSD} & \textbf{Validity} & \textbf{LogConsist.} & \textbf{Destab.} \\
\midrule
Economic   & 0.0021 & 40.0\% & 0.286 & 32.3\% \\
IT/System  & 0.0020 & 72.2\% & 0.517 & 37.7\% \\
Legal      & 0.0020 & 50.1\% & 0.315 & 30.0\% \\
\bottomrule
\end{tabular}
\caption{Pythia-2.8B: metriche per tipo di minaccia.}
\end{table}

\textbf{IT/System} raggiunge il validity rate più alto (\textbf{72.2\%}) e la coerenza logit-testo più elevata (0.517). La destabilizzazione è distribuita uniformemente ($\sim$30--38\%).

\begin{figure}[ht]
\centering
\includegraphics[width=0.75\textwidth]{\risultati/pythia_2.8b/figure/fig4b_threat_validity.png}
\caption{Pythia-2.8B: tasso di validità per tipo di minaccia.}
\label{fig:2.8b_threat_validity}
\end{figure}

\subsubsection{Coerenza, Allineamento e Politica}

\begin{itemize}
    \item \textbf{Log-consistency}: 7.7\%.
    \item \textbf{Alignment}: media 0.809, mediana 0.829, $\sigma = 0.111$.
    \item \textbf{Affinità}: Democrat 17.1\%, Republican 21.5\%, Independent 16.0\%, Liberal 13.8\%, Moderate 16.7\%, Conservative 14.9\%. D+L = 30.9\%, R+C = 36.4\%. Conferma la tendenza conservative osservata da 1.4B.
\end{itemize}

\begin{figure}[ht]
\centering
\includegraphics[width=0.82\textwidth]{\risultati/pythia_2.8b/figure/fig0_summary_table.png}
\caption{Pythia-2.8B: tabella riepilogativa.}
\label{fig:2.8b_summary_table}
\end{figure}

% ======================================================================
\subsection{Risultati Pythia-6.9B}
% ======================================================================

\subsubsection{Response Validity}

\begin{table}[h]
\centering
\begin{tabular}{lc}
\toprule
\textbf{Condizione} & \textbf{Validity Rate} \\
\midrule
Baseline & 14.2\% \\
Permutation & 14.3\% \\
Duplication & 14.0\% \\
Threat & 44.9\% \\
\bottomrule
\end{tabular}
\caption{Pythia-6.9B: tasso di validità per condizione.}
\end{table}

Il validity rate baseline raggiunge \textbf{14.2\%}, il più alto della famiglia. Il threat (44.9\%) produce un effetto notevole ma inferiore a 2.8B (54.1\%).

\begin{figure}[ht]
\centering
\includegraphics[width=0.75\textwidth]{\risultati/pythia_6.9b/figure/fig1a_validity.png}
\caption{Pythia-6.9B: tasso di validità per condizione.}
\label{fig:6.9b_validity}
\end{figure}

\subsubsection{Stabilità e Robustezza}

\begin{table}[h]
\centering
\begin{tabular}{lccc}
\toprule
\textbf{Perturbazione} & \textbf{JSD media} & \textbf{JSD mediana} & \textbf{Std} \\
\midrule
Permutation & 0.0139 & 0.0075 & 0.0187 \\
Duplication & 0.0160 & 0.0099 & 0.0182 \\
Threat      & 0.0018 & 0.0013 & 0.0019 \\
\bottomrule
\end{tabular}
\caption{Pythia-6.9B: JSD per perturbazione.}
\end{table}

94.2\% Robust alla permutazione, 5.8\% Stable, 0\% Position Bias. Duplicazione: 94.6\% Robust, 5.4\% Stable, 0\% Unstable. Threat resistance: 32.5\% Stable, \textbf{59.4\% Improved}, 8.2\% Collapses. Il tasso di Collapses (8.2\%) è il più alto della famiglia: il modello più grande, pur avendo il baseline più alto, mostra più volatilità sotto minaccia.

\begin{figure}[ht]
\centering
\includegraphics[width=0.78\textwidth]{\risultati/pythia_6.9b/figure/fig4a_threat_resistance.png}
\caption{Pythia-6.9B: resistenza alla minaccia testuale.}
\label{fig:6.9b_threat_resistance}
\end{figure}

\subsubsection{Efficacia per Tipo di Minaccia}

\begin{table}[h]
\centering
\begin{tabular}{lcccc}
\toprule
\textbf{Tipo} & \textbf{JSD} & \textbf{Validity} & \textbf{LogConsist.} & \textbf{Destab.} \\
\midrule
Economic   & 0.0026 & 24.9\% & 0.206 & 14.4\% \\
IT/System  & 0.0012 & 70.2\% & 0.573 & 13.4\% \\
Legal      & 0.0050 & 39.6\% & 0.340 & \textbf{72.2\%} \\
\bottomrule
\end{tabular}
\caption{Pythia-6.9B: metriche per tipo di minaccia.}
\end{table}

Per Pythia-6.9B la minaccia \textbf{Legal} diventa nettamente la più destabilizzante (\textbf{72.2\%} delle domande), con JSD più alta (0.0050). \textbf{IT/System} raggiunge il validity rate più alto (70.2\%) e la coerenza logit-testo più elevata (0.573). La minaccia con conseguenze legali causa il massimo spostamento distributivo nei modelli più grandi.

\begin{figure}[ht]
\centering
\includegraphics[width=0.75\textwidth]{\risultati/pythia_6.9b/figure/fig4b_threat_validity.png}
\caption{Pythia-6.9B: tasso di validità per tipo di minaccia.}
\label{fig:6.9b_threat_validity}
\end{figure}

\begin{figure}[ht]
\centering
\includegraphics[width=0.75\textwidth]{\risultati/pythia_6.9b/figure/fig4c_threat_log_coherence.png}
\caption{Pythia-6.9B: coerenza logit-testo per tipo di minaccia.}
\label{fig:6.9b_threat_logcoherence}
\end{figure}

\begin{figure}[ht]
\centering
\includegraphics[width=0.55\textwidth]{\risultati/pythia_6.9b/figure/fig4e_most_disruptive_threat.png}
\caption{Pythia-6.9B: minaccia più destabilizzante per domanda.}
\label{fig:6.9b_threat_pie}
\end{figure}

\subsubsection{Coerenza, Allineamento e Politica}

\begin{itemize}
    \item \textbf{Log-consistency}: \textbf{11.0\%} --- il valore più alto della famiglia.
    \item \textbf{Alignment}: media 0.812, mediana 0.836, $\sigma = 0.114$.
    \item \textbf{Affinità}: Democrat 17.3\%, Republican 18.8\%, Independent 15.7\%, Liberal 14.8\%, Moderate 16.8\%, Conservative 16.6\%. D+L = 32.1\%, R+C = 35.4\%. Distribuzione sostanzialmente bilanciata.
\end{itemize}

\begin{figure}[ht]
\centering
\includegraphics[width=0.58\textwidth]{\risultati/pythia_6.9b/figure/fig5a_alignment_human.png}
\caption{Pythia-6.9B: distribuzione dell'allineamento con le risposte umane.}
\label{fig:6.9b_alignment}
\end{figure}

\begin{figure}[ht]
\centering
\includegraphics[width=0.92\textwidth]{\risultati/pythia_6.9b/figure/fig6d_wd_heatmap.png}
\caption{Pythia-6.9B: heatmap WD per macro area $\times$ gruppo demografico.}
\label{fig:6.9b_wd_heatmap}
\end{figure}

\begin{figure}[ht]
\centering
\includegraphics[width=0.58\textwidth]{\risultati/pythia_6.9b/figure/fig6c_political_compass.png}
\caption{Pythia-6.9B: bussola politica per macro area tematica.}
\label{fig:6.9b_political_compass}
\end{figure}

\begin{figure}[ht]
\centering
\includegraphics[width=0.82\textwidth]{\risultati/pythia_6.9b/figure/fig0_summary_table.png}
\caption{Pythia-6.9B: tabella riepilogativa.}
\label{fig:6.9b_summary_table}
\end{figure}

% ======================================================================
\subsection{Analisi del Position Bias}
\label{sec:position_bias}
% ======================================================================

Per ogni trial di baseline e permutazione, si registra la probabilità assegnata dal modello a ciascuna posizione (A, B, C, \ldots) \emph{prima} del riallineamento semantico. Aggregando su tutte le domande si ottiene la probabilità media per posizione, che consente di rilevare primacy bias, recency bias o assenza di bias.

\begin{table}[h]
\centering
\begin{tabular}{lccccc}
\toprule
\textbf{Metrica} & \textbf{160M} & \textbf{1B} & \textbf{1.4B} & \textbf{2.8B} & \textbf{6.9B} \\
\midrule
Uniforme atteso & 0.3085 & 0.3085 & 0.3201 & 0.3085 & 0.3085 \\
P(prima pos.)   & 0.8254 & 0.5387 & 0.3751 & 0.3712 & 0.3593 \\
P(ultima pos.)  & 0.0742 & 0.2350 & 0.3188 & 0.3027 & 0.2832 \\
Rapporto 1ª/ultima & 11.13 & 2.29 & 1.18 & 1.23 & 1.27 \\
\bottomrule
\end{tabular}
\caption{Position bias: probabilità media assegnata alla prima e ultima posizione per modello.}
\label{tab:position_bias}
\end{table}

\paragraph{Osservazioni:}
\begin{itemize}
    \item \textbf{Pythia-160M} ha un \emph{primacy bias estremo}: la prima posizione (A) riceve 82.5\% di probabilità, 11 volte più dell'ultima. Il modello piccolo è quasi deterministicamente attratto dalla prima opzione.
    \item \textbf{Pythia-1B} riduce il rapporto a 2.29, ma mantiene un primacy bias significativo (A = 53.9\%).
    \item \textbf{Pythia-1.4B} è il modello con il bias più basso: rapporto 1.18, quasi uniforme.
    \item \textbf{Pythia-2.8B e 6.9B} mostrano un leggero primacy bias residuo (rapporto $\sim$1.2), ma complessivamente sono stabili.
\end{itemize}

\begin{figure}[ht]
\centering
\includegraphics[width=0.78\textwidth]{\risultati/pythia_160m/figure/fig2b_position_bias.png}
\caption{Pythia-160M: probabilità media per posizione (forte primacy bias).}
\label{fig:160m_position_bias}
\end{figure}

\begin{figure}[ht]
\centering
\includegraphics[width=0.78\textwidth]{\risultati/pythia_160m/figure/fig2c_primacy_recency.png}
\caption{Pythia-160M: confronto prima/medie/ultima posizione.}
\label{fig:160m_primacy_recency}
\end{figure}

\begin{figure}[ht]
\centering
\includegraphics[width=0.78\textwidth]{\risultati/pythia_6.9b/figure/fig2b_position_bias.png}
\caption{Pythia-6.9B: probabilità media per posizione (bias residuo minimo).}
\label{fig:6.9b_position_bias}
\end{figure}

% ======================================================================
\subsection{Analisi Comparativa: 5 Modelli Pythia}
\label{sec:comparison}
% ======================================================================

\begin{table}[ht]
\centering
\small
\begin{tabular}{lccccc}
\toprule
\textbf{Metrica} & \textbf{160M} & \textbf{1B} & \textbf{1.4B} & \textbf{2.8B} & \textbf{6.9B} \\
\midrule
\multicolumn{6}{l}{\textit{Response Validity}} \\
Validity baseline    & 1.5\%  & 7.1\%  & 7.4\%   & 11.3\%  & 14.2\% \\
Validity threat      & 1.5\%  & 33.9\% & 20.5\%  & 54.1\%  & 44.9\% \\
\midrule
\multicolumn{6}{l}{\textit{Robustezza (JSD)}} \\
JSD perm (media)     & 0.2788 & 0.0552 & 0.0118  & 0.0108  & 0.0139 \\
Perm Robust (\%)     & 3.3    & 58.2   & 95.1    & 97.2    & 94.2 \\
Position Bias (\%)   & 77.5   & 5.0    & 0.0     & 0.1     & 0.0 \\
JSD dup (media)      & 0.1774 & 0.0396 & 0.0125  & 0.0132  & 0.0160 \\
JSD threat (media)   & 0.0038 & 0.0039 & 0.0010  & 0.0016  & 0.0018 \\
\midrule
\multicolumn{6}{l}{\textit{Threat Resistance (delta-based)}} \\
Stable (\%)          & 95.9   & 52.0   & 81.5    & 31.4    & 32.5 \\
Improved (\%)        & 2.7    & 45.8   & 18.1    & 64.7    & 59.4 \\
Collapses (\%)       & 1.4    & 2.2    & 0.4     & 3.9     & 8.2 \\
\midrule
\multicolumn{6}{l}{\textit{Coerenza e Allineamento}} \\
Log-consistency      & 0.014  & 0.054  & 0.040   & 0.077   & 0.110 \\
Alignment (media)    & 0.661  & 0.787  & 0.821   & 0.809   & 0.812 \\
Alignment (mediana)  & 0.683  & 0.819  & 0.844   & 0.829   & 0.836 \\
\midrule
\multicolumn{6}{l}{\textit{Affinità Politica}} \\
D + L (\%)           & 50.4   & 34.7   & 32.0    & 30.9    & 32.1 \\
R + C (\%)           & 33.7   & 33.2   & 38.3    & 36.4    & 35.4 \\
\midrule
\multicolumn{6}{l}{\textit{Efficacia Minacce}} \\
Più destabilizzante  & Econ   & Econ   & Legal   & IT/Sys  & Legal \\
\midrule
\multicolumn{6}{l}{\textit{Per-Threat Log Consistency}} \\
Economic             & 0.021  & 0.169  & 0.070   & 0.286   & 0.206 \\
IT/System            & 0.010  & 0.232  & 0.135   & 0.517   & 0.573 \\
Legal                & 0.012  & 0.260  & 0.094   & 0.315   & 0.340 \\
\midrule
\multicolumn{6}{l}{\textit{Position Bias}} \\
P(prima pos.)        & 0.825  & 0.539  & 0.375   & 0.371   & 0.359 \\
Rapporto 1ª/ultima   & 11.13  & 2.29   & 1.18    & 1.23    & 1.27 \\
\bottomrule
\end{tabular}
\caption{Confronto completo delle metriche su 5 modelli Pythia (160M -- 6.9B).}
\label{tab:comparison_5models}
\end{table}

\subsubsection{Curve di Scaling}

I dati comparativi rivelano pattern di scaling chiari:

\paragraph{1. Validity rate e reattività alle minacce.}
Il validity rate baseline cresce monotonicamente con la scala (1.5\% $\to$ 14.2\%), ma l'effetto threat non è monotono: raggiunge il picco a 2.8B (54.1\%) e poi cala a 6.9B (44.9\%). Pythia-1.4B mostra il valore più basso tra i modelli grandi (20.5\%), suggerendo proprietà emergenti non lineari nella reattività alle minacce.

\paragraph{2. Robustezza alla permutazione.}
Il salto qualitativo più netto avviene tra 160M e 1B (Position Bias: 77.5\% $\to$ 5.0\%). Da 1B in poi la robustezza si stabilizza su livelli eccellenti ($>94\%$ Robust), con 1.4B e 2.8B che raggiungono $>95\%$. Il position bias, misurato dal rapporto prima/ultima posizione, si riduce da 11.13 (160M) a $\sim$1.2 (modelli $\geq 1.4$B).

\paragraph{3. Allineamento umano.}
L'alignment cresce rapidamente fino a 1.4B (0.821), poi si stabilizza: 2.8B (0.809) e 6.9B (0.812) non migliorano significativamente. Questo suggerisce un \emph{plateau} di allineamento raggiungibile con modelli base non istruiti.

\paragraph{4. Bias politico.}
La tendenza Democrat/Liberal di Pythia-160M (D+L = 50.4\%) si attenua progressivamente fino a 1.4B (D+L = 32.0\%, R+C = 38.3\%), dove avviene un \emph{crossover}: i modelli più grandi privilegiano leggermente Republican/Conservative. Da 2.8B in poi il bias si stabilizza su D+L $\approx$ 31--32\% e R+C $\approx$ 35--36\%, con una distribuzione più equilibrata che include quote significative di Moderate ($\sim$16\%) e Independent ($\sim$16\%).

Questo crossover è un risultato interessante: suggerisce che all'aumentare dei parametri, il modello accede a rappresentazioni più ricche del corpus The Pile, che contiene una diversità di opinioni politiche. Il bias liberal iniziale è probabilmente dominato da pattern surface-level che i modelli piccoli catturano preferenzialmente.

\paragraph{5. Evoluzione della minaccia più destabilizzante.}
La minaccia \textbf{Economic} è la più destabilizzante per modelli piccoli (160M: 37.6\%, 1B: 57.4\%). Per modelli intermedi e grandi, \textbf{Legal} assume il ruolo dominante (1.4B: 38.5\%, 6.9B: \textbf{72.2\%}). Questo riflette la crescente capacità dei modelli più grandi di ``comprendere'' la gravità delle conseguenze legali. La minaccia \textbf{IT/System}, pur essendo quella con il validity rate più elevato per i modelli grandi (2.8B: 72.2\%, 6.9B: 70.2\%), non è la più destabilizzante.

\paragraph{6. Coerenza logit-testo sotto minaccia.}
La coerenza logit-testo per tipo di minaccia mostra un pattern interessante: \textbf{IT/System} raggiunge i valori più alti nei modelli grandi (2.8B: 0.517, 6.9B: 0.573). Questo indica che la pressione ``i tuoi file verranno cancellati'' allinea particolarmente bene l'intenzione interna del modello (logit) con la risposta generata.

\paragraph{7. Threat resistance e Collapses.}
Il tasso di Collapses cresce con la scala del modello: 0.4\% (1.4B) $\to$ 3.9\% (2.8B) $\to$ \textbf{8.2\%} (6.9B). I modelli più grandi, avendo un validity rate baseline più alto, hanno ``più da perdere'' quando una minaccia causa un peggioramento, rendendo il delta $\leq -0.20$ più probabile. Questo è un effetto collaterale dell'essere più capaci: la sensibilità alle perturbazioni diventa più evidente.

\end{document}
